\documentclass[10pt,a4paper]{amsart}
\usepackage[margin=19mm]{geometry}
\usepackage{amsmath,amssymb,mathtools}
\usepackage[scr=rsfs,cal=euler]{mathalfa}
\usepackage{xfrac}
\usepackage[unicode,hidelinks,bookmarks=false]{hyperref}
\newcommand{\ie}{i.e.\ }
\newcommand{\cf}{cf.\ }
\newcommand{\resp}{respectively\ }
\newcommand{\Subsection}{\subsection}
\providecommand{\nobreakdash}{\nobreak\hbox{-}}
\setlength{\emergencystretch}{2em}
\multlinegap0pt
\usepackage{bm}
\usepackage{bbm}
\usepackage{dsfont}
\usepackage{xspace}
\usepackage{cancel}
\usepackage{mathtools}
\usepackage{amsthm}
\makeatletter
\@ifundefined{thm}{\newtheorem{thm}{Thm}}{}
\@ifundefined{proposition}{\newtheorem{proposition}[thm]{Proposition}}{}
\makeatother
\def\I{\mathrm{i}}
\title[Two dimensional potential theory with a view towards vortex ]{Two dimensional potential theory with a view towards vortex motion: Energy, capacity and Green functions}
\author{Bj\"orn Gustafsson}
\date{Source: arXiv:2405.19215v3 (CC BY 4.0)}
\begin{document}
\maketitle

\noindent\emph{Source and licence.} Two dimensional potential theory with a view towards vortex motion: Energy, capacity and Green functions. Version arXiv:2405.19215v3 (CC BY 4.0), distributed under the Creative Commons Attribution 4.0 International licence, as recorded in the arXiv licence metadata for this exact version and shown on the abstract page https://arxiv.org/abs/2405.19215v3. Full rights notice: licenses/arxiv-2405.19215-CC-BY-4.0.txt.\par\smallskip

\noindent\textit{Editorial scope.} Counted unit: the Proposition whose source label is \texttt{prop:GUU} in the article (source0.tex lines 2533--2540, source label \texttt{prop:GUU}), delivered together with the article's own complete original proof of it (source0.tex lines 2542--2572, one \texttt{proof} environment of 31 lines). No mathematics is written, repaired or re-derived: statement and proof are the article's own text line for line. Required context retained uncounted: omegaeta (lines 2249--2253); periods (lines 2309--2318); PRRQ (lines 2321--2330). Every definition, display and label of those lines is retained unchanged. Omitted from this package and not redistributed: 1--2248, 2254--2308, 2319--2320, 2331--2532, 2541--2541, 2573--3955. Nothing inside a retained excerpt is omitted or altered: every retained body is delivered exactly as the article's lines are written, so no omission receipt of any kind accompanies this package. Citations: the retained excerpts carry no citation command at all, so no bibliography entry of the article is transcribed and no reference is redistributed. Reference adaptation: none was needed and none was made; every reference inside the delivered unit resolves inside it, so no unresolved reference or citation is produced. Printed slips kept literal: no printed word, symbol, hypothesis or number of the counted statement or of its proof was altered. Layout and class: the article's own class is \texttt{article} with packages including palatino, amssymb, amsthm, amsmath, amsfonts, epsfig, helvet, courier, type1cm, makeidx, graphicx, multicol, footmisc; the wrapper is the lane's pinned \texttt{amsart} editorial wrapper at 10pt on A4, which restates the theorem-like environments the retained text needs and adds the article's own macro definitions that the retained text needs (\texttt{\textbackslash I}), with extra wrapper packages (bm, bbm, dsfont, xspace, cancel, mathtools). The delivered numbering is the unit's own. Assets: no retained span contains a figure environment, a graphics-inclusion command, a PDF-page inclusion, a table or a TikZ picture; no image file is copied into this package, the package declares no asset dependency and no image file is redistributed with it. Licence: version arXiv:2405.19215v3 of this article is distributed under the Creative Commons Attribution 4.0 International licence, as recorded in the arXiv licence metadata for this exact version and shown on the abstract page https://arxiv.org/abs/2405.19215v3. A full rights notice is delivered at licenses/arxiv-2405.19215-CC-BY-4.0.txt. Characters: every retained excerpt is pure ASCII, so the delivered engine, XeLaTeX with its default Latin Modern text font, reports no missing character.
\begin{align}\label{omegaeta}
\omega_\gamma(a)=\eta_\gamma(a)+\I*\eta_\gamma(a)
&=\I \oint_\gamma L(z,a)dzda\\ 
&=-\I \oint_\gamma \overline{K(z,a)dzd\bar{a}}.
\end{align}

\begin{equation}\label{periods}
\left( \begin{array}{cc}
-\oint_{\alpha} \eta_{\beta}^T& -\oint_{\beta} \eta_{\beta}^T\\
\oint_{\alpha} \eta_{\alpha}^T & \oint_{\beta} \eta_{\alpha}^T 
\end{array} \right)
=\left( \begin{array}{cc}
                                            I & 0  \\
                                             0 & I  \\
\end{array} \right).
\end{equation}

\begin{equation}\label{PRRQ}
\left( \begin{array}{cc}
                                            P & R  \\
                                              R^T & Q  \\
\end{array} \right)=
\left( \begin{array}{cc} 
                                             - \oint_{\beta}*\eta_{\beta}^T & \oint_{\beta}*\eta_{\alpha}^T  \\
                                               \oint_{\alpha}*\eta_{\beta}^T & ( -\oint_{\alpha}*\eta_{\alpha}^T  \\
\end{array} \right).
\end{equation}

\begin{proposition}\label{prop:GUU} 
In terms of the matrices $P$, $Q$ we have 
\begin{align}\label{KPQ} 
K(z,a)dzd\bar{a}
&=\sum_{k,j=1}^\texttt{g} (P^{-1})_{kj}\omega_{\beta_k}(z)\overline{\omega_{\beta_j}(a)}\\
&=\sum_{k,j=1}^\texttt{g} (Q^{-1})_{kj}\omega_{\alpha_k}(z)\overline{\omega_{\alpha_j}(a)}.
\end{align}
\end{proposition}

\begin{proof}
Since $K(z,a)dzda$ is holomorphic in $z$,
anti-holomorphic in $a$, there must be expansions of the form in the lemma, for some choice of coefficients.
It only remains to check that the coefficients in the proposition are the right ones. 

On keeping $a$ fixed this is confirmed by showing that either all integrals
$\oint_{\alpha_j}dz$, or all integrals $\oint_{\beta_j}dz$, 
come out the same when applied to the left and right members above. 
If one of these sets of integrals come out correctly, then the other set must too, because an everywhere holomorphic 
differential on $M$ is determined by its periods around any one of these sets of cycles.

The check is indeed straight-forward. Combining (\ref{periods}) and (\ref{PRRQ}) we have  
$$
\oint_\alpha \omega_{\alpha}^T =-\I Q, \quad \oint_\alpha \omega_\beta^T =\I R^T -I,
$$
$$
\oint_\beta \omega_\alpha^T =\I R+I, \quad \oint_\beta \omega_\beta^T =-\I P.
$$
This gives, in view of (\ref{omegaeta}) and with integration with respect to $z$, 
$$
\oint_{\alpha}\omega_\alpha(z)^T Q^{-1}\overline{\omega_\alpha(a)} 
=-\I Q^T Q^{-1}\overline{\omega_\alpha(a)}=-\I\overline{\omega_\alpha(a)}  
= \oint_{\alpha}K(z,a)dzd\bar{a}, 
$$
$$ 
\oint_{\beta}\omega_{\beta}^T(z)P^{-1}\overline{\omega_\beta(a)} 
=-\I P^T P^{-1}\overline{\omega_\beta(a)}=-\I\overline{\omega_\beta(a)}  
=\oint_{\beta}K(z,a)dzd\bar{a},
$$ 
as required.
\end{proof}

\end{document}
